\chapter{Формальная постановка задачи}
\label{cha:Formal}

В данной главе вводятся математические обозначения, описываются форматы входных и выходных данных, а также формализуются метрики качества укладки.

\section{Математическая модель}

Паллета моделируется прямоугольным параллелепипедом с основанием $L \times W$ и нормировочной высотой $H$. В экспериментах используется стандартная паллета EUR $1200 \times 800$~мм с нормировочной высотой $H = 2200$~мм.

Дан набор из $K$ типов коробок. Коробка типа $i$ характеризуется размерами $(l_i, w_i, h_i)$, массой $m_i$ и количеством экземпляров $q_i$. Общее число коробок составляет $N = \sum_{i=1}^{K} q_i$.

Решением задачи является отображение каждого экземпляра коробки в позицию и ориентацию на паллете. Позиция $j$-го экземпляра задаётся координатами двух противоположных вершин: нижней левой $(x_j, y_j, z_j)$ и верхней правой $(X_j, Y_j, Z_j)$, где $X_j - x_j$ и $Y_j - y_j$ --- длина и ширина коробки в выбранной ориентации, а $Z_j - z_j = h_i$ --- высота.

Допускается единственный тип вращения --- поворот на $90^{\circ}$ вокруг вертикальной оси, при котором $l_i$ и $w_i$ меняются местами. Переворот коробки (перестановка высоты с длиной или шириной) запрещён.

Решение является допустимым, если выполнены следующие условия:
\begin{enumerate}
  \item \textbf{Вложенность:} $0 \leq x_j$, $X_j \leq L$, $0 \leq y_j$, $Y_j \leq W$ для всех $j$;
  \item \textbf{Непересечение:} для любых двух коробок $j \neq k$ их объёмные проекции не пересекаются;
  \item \textbf{Опора:} $z_j = 0$ (коробка стоит на паллете) либо $z_j$ совпадает с верхней гранью $Z_k$ хотя бы одной коробки $k$, проекция которой перекрывает центр нижней грани $j$-й коробки в плане $XY$: $\bigl(x_j + \tfrac{X_j - x_j}{2},\; y_j + \tfrac{Y_j - y_j}{2}\bigr)$.
\end{enumerate}

\section{Формат входных данных}
\label{sec:input_format}

Входные данные представляют собой CSV-файл, каждая строка которого описывает один тип коробок:

\begin{center}
\texttt{SKU, Quantity, Length, Width, Height, Weight, Strength, Aisle, Caustic}
\end{center}

\noindent где:
\begin{itemize}
  \item \textbf{SKU} --- уникальный идентификатор товарной позиции;
  \item \textbf{Quantity} --- количество экземпляров данного типа;
  \item \textbf{Length}, \textbf{Width}, \textbf{Height} --- размеры коробки в миллиметрах;
  \item \textbf{Weight} --- масса коробки в граммах;
  \item \textbf{Strength} --- прочность (в текущей версии не используется);
  \item \textbf{Aisle} --- номер секции магазина (категория товара);
  \item \textbf{Caustic} --- признак едкого содержимого (в текущей версии не используется).
\end{itemize}

\section{Формат выходных данных}
\label{sec:output_format}

Результат укладки описывается CSV-файлом, каждая строка которого задаёт позицию одной коробки:

\begin{center}
\texttt{Pallet, SKU, x, y, z, X, Y, Z}
\end{center}

\noindent где \texttt{Pallet} --- номер паллеты (от 0), \texttt{SKU} --- идентификатор типа коробки, $(x, y, z)$ --- координаты нижнего левого угла, $(X, Y, Z)$ --- координаты верхнего правого угла. Порядок строк определяет последовательность укладки.

\section{Метрики качества укладки}

Для оценки решения используется набор метрик, каждая из которых отражает один из аспектов качества: плотность заполнения, устойчивость и равномерность загрузки. Все метрики принимают значения в диапазоне $[0, 1]$ и объединяются в скалярную целевую функцию для оптимизации.

\section{Коэффициент заполнения (percolation)}
\label{sec:metric_perc}

Коэффициент заполнения определяется как отношение суммарного объёма уложенных коробок к объёму паллеты:
\[
\text{percolation} = \frac{\sum_{j=1}^{N} V_j}{L \times W \times H_{\max}},
\]
где $V_j = (X_j - x_j)(Y_j - y_j)(Z_j - z_j)$ --- объём $j$-й коробки, а $H_{\max} = \max_j Z_j$ --- фактическая высота штабеля. В случае нескольких паллет перколяция вычисляется как совокупная по всем паллетам: сумма объёмов коробок делится на сумму объёмов паллет.

Данная метрика принимает значения в диапазоне $(0, 1]$. Чем выше перколяция, тем эффективнее используется объём паллеты.

\section{Минимальный коэффициент опоры (min\_support\_ratio)}
\label{sec:metric_support}

Устойчивость укладки оценивается через минимальный коэффициент опоры --- наименьшую среди всех коробок долю площади нижней грани, опирающуюся на нижележащий слой.

Для коробки $j$ с нижней гранью площадью $A_j = (X_j - x_j)(Y_j - y_j)$ коэффициент опоры определяется как:
\[
\text{support\_ratio}_j = \frac{A_j^{\text{supported}}}{A_j},
\]
где $A_j^{\text{supported}}$ --- площадь области нижней грани коробки~$j$, которая опирается на верхнюю грань нижележащих коробок (для коробок на поверхности паллеты $\text{support\_ratio}_j = 1$). Метрика всей укладки:
\[
\text{min\_support\_ratio} = \min_{j=1}^{N}\, \text{support\_ratio}_j.
\]

\paragraph{Обоснование выбора метрики.}
В промышленной практике для оценки устойчивости используется ряд подходов: ограничение на долю площади опоры~\cite{Bischoff1995}, проверка нахождения центра тяжести внутри полигона опоры~\cite{Calzavara2021}, stacking-tree с физическим моделированием~\cite{Zhao2021}. 

Метрика min\_support\_ratio выбрана по нескольким причинам:
\begin{itemize}
  \item \textbf{Вычислительная эффективность.} Расчёт площади опоры реализуется через структуру HeightHandler за время $O(n)$ для каждой коробки, что не увеличивает асимптотику основного алгоритма.
  \item \textbf{Дифференцируемость по решению.} Метрика непрерывно зависит от расположения коробок, что позволяет использовать её в составе целевой функции метаэвристических алгоритмов с градиентоподобной оптимизацией.
  \item \textbf{Интерпретируемость.} Значение 0,25 означает, что хотя бы одна коробка опирается лишь четвертью основания, что легко визуализировать и оценить.
  \item \textbf{Чувствительность к наихудшему случаю.} Операция минимума гарантирует, что ни одна коробка не <<повисает>> --- в отличие от усреднённых метрик, которые допускают наличие полностью неопёртых коробок при высоком среднем значении.
\end{itemize}

Дополнительно в алгоритме применяется ограничение на поддержку центра масс (см.~раздел~\ref{sec:stability_constraint}): коробка может быть уложена только в позицию, где геометрический центр её основания проецируется на опорную поверхность. Это ограничение пассивно поднимает min\_support\_ratio до значений не менее 0,25.

\section{Баланс высот (height\_balance)}
\label{sec:metric_hb}

При укладке на несколько паллет важна равномерность загрузки. Метрика баланса высот определяется через коэффициент вариации высот штабелей:
\[
\text{height\_balance} = \max\!\left(0,\; 1 - \frac{\sigma(H_1, \ldots, H_P)}{\mu(H_1, \ldots, H_P)}\right),
\]
где $H_i$ --- высота штабеля на $i$-й паллете, $\mu$ --- среднее, $\sigma$ --- стандартное отклонение. При $P = 1$ метрика равна~1.

Значение 1 соответствует идеально равномерной загрузке, а значения ниже 0,8 указывают на существенный дисбаланс, при котором одна паллета нагружена значительно больше других.

\section{Нормализованная координата центра масс}
\label{sec:metric_com}

Для оценки распределения веса по высоте используется нормализованная $z$-координата центра масс:
\[
\text{CoM}_z^{\text{rel}} = \frac{\sum_{j=1}^{N} m_j \cdot \frac{z_j + Z_j}{2}}{\left(\sum_{j=1}^{N} m_j\right) \cdot H_{\max}}.
\]

Низкие значения $\text{CoM}_z^{\text{rel}}$ (ближе к 0) соответствуют укладке, в которой тяжёлые коробки расположены внизу, что повышает физическую устойчивость штабеля при транспортировке. Метрика входит в целевую функцию с отрицательным знаком: оптимизируется $(1 - \text{CoM}_z^{\text{rel}})$.

\section{Целевая функция}

Все метрики объединяются в скалярную целевую функцию:
\[
\begin{split}
\text{score} = {}& w_1 \cdot \text{percolation}
    + w_2 \cdot \text{min\_support\_ratio} \\
&{}+ w_3 \cdot (1 - \text{CoM}_z^{\text{rel}})
    + w_4 \cdot \text{height\_balance},
\end{split}
\]
где $w_1, \ldots, w_4$ --- настраиваемые неотрицательные веса. Для задачи с одной паллетой типичными являются значения $w_1 = 1$, $w_2 = 2$, $w_3 = w_4 = 0$, подобранные экспериментально (см.~раздел~\ref{sec:evyukh_experiments}).

\section{Набор данных}

Экспериментальное тестирование проводилось на наборе из 436 реальных задач, предоставленных компанией-оператором складской логистики. Каждая задача содержит от 3 до 120 типов коробок (от 5 до 400 экземпляров) и охватывает широкий диапазон размеров: от $50 \times 50 \times 30$~мм до $600 \times 400 \times 400$~мм. Масса отдельных коробок варьируется от 100 г до 25 кг.

Для мультипаллетного режима из одиночных задач формируются составные тесты путём случайного объединения нескольких заказов в одну группу доставки. Количество паллет в составных тестах варьируется от 2 до 6.
